\documentclass[11pt,a4paper]{article}
\usepackage[margin=2cm]{geometry}
\usepackage{enumitem}
\usepackage[dvipsnames]{xcolor}
\usepackage{hyperref}
\usepackage{titlesec}
\usepackage{mdframed}
\hypersetup{colorlinks=true,linkcolor=NavyBlue}
\setlength{\parindent}{0pt}
\setlength{\parskip}{5pt}
\setlist{itemsep=1pt,topsep=2pt}
\titleformat{\section}{\Large\bfseries\color{NavyBlue}}{}{0em}{}
\newcounter{slide}
\newcommand{\slide}[2]{\par\bigskip\stepcounter{slide}%
  {\color{NavyBlue}\rule{\linewidth}{0.4pt}}\par
  \textbf{\large Slide \arabic{slide}: #1}\hfill{\small\color{gray}#2}\par}
\newcommand{\act}[1]{\par{\small\color{OliveGreen}\textit{[#1]}}\par}
\newmdenv[linecolor=BrickRed,linewidth=0.8pt,backgroundcolor=BrickRed!4,innertopmargin=4pt,innerbottommargin=4pt]{warnbox}
\newmdenv[linecolor=NavyBlue,linewidth=0.8pt,backgroundcolor=NavyBlue!4,innertopmargin=4pt,innerbottommargin=4pt]{keybox}

\title{\textbf{Presentation script}\\[2pt]\large Short Solutions for Large Rubik's Cubes and Squares}
\author{Slide-by-slide spoken script for \texttt{Rubik\_Short\_Solutions.pptx} (35 slides)}
\date{18 November 2026}

\begin{document}
\maketitle

\begin{keybox}
\textbf{How to use this.} Plain text is what you say, in simple words. \textcolor{OliveGreen}{\textit{[Green brackets]}} are
what to do or point at. Times are a guide: read in full, the script runs about \textbf{30 minutes}. For a 15--20 minute slot, skip the slides
marked \textbf{(optional)} and say only the \emph{first paragraph} on each slide; the main points are always there. Speak slowly; pause after every big number.
\end{keybox}

\textbf{Suggested split for six speakers} (change as you like):
\begin{itemize}
\item Speaker 1: slides 1--5 (problem and questions)
\item Speaker 2: slides 6--11 (how the puzzles work, the engine, the cast of solvers)
\item Speaker 3: slides 12--16 (the reproductions and OneNet)
\item Speaker 4: slides 17--22 (NetGain, SymPort, the bound)
\item Speaker 5: slides 23--31 (results)
\item Speaker 6: slides 32--35 (discussion, limits, future work, close)
\end{itemize}
Hand-over line: ``Now [name] will explain \ldots''

\begin{warnbox}
\textbf{Three slide numbers that disagree with the paper. Say the paper's number.}
\begin{itemize}
\item Slide 16 shows ``$r\ge0.70$''. The paper's numbers are \textbf{0.66, 0.65 and 0.53} at sizes 10, 14 and 20. The script below
uses the paper's numbers. Better still, fix the slide.
\item Slides 24 and 32 say one run of KMCoders' program is ``1.07--1.18$\times$'' / ``7--18\% longer''. The data say
\textbf{1.05--1.18$\times$} (5--18\%).
\item If asked about ``21\% at $N=10$'' (abstract): that is from 2 random cubes; the 20-cube test gives \textbf{17\%}.
\end{itemize}
\end{warnbox}

\newpage
%% ====================================================================
\section{Part 1: The problem}

\slide{Title}{30 s}
Good morning. We are [names]. Our project is about solving big Rubik's cubes, and their flat cousin, the Rubik's Square, in
as few moves as possible.

In one sentence: we built one engine that handles cubes and squares of every size, ran eight different solvers on exactly the
same puzzles, and our main method, \textbf{SymPort}, beats the best published competition solutions on average for cubes from
4 by 4 up to 7 by 7. We will also show honestly where and why it stops winning.

\slide{Problem statement}{1 min}
The task is simple to say: you get a scrambled cube of any size; give back a short list of moves that solves it.

It is hard for three reasons. \act{point at the three boxes, left to right}
First, the number of possible states explodes as the cube gets bigger. A 20 by 20 by 20 cube has 2,400 stickers.

Second, finding the \emph{shortest} solution is NP-complete. That means no one knows a fast method that is guaranteed to
find the best answer, and most experts think none exists. Finding \emph{some} solution is easy; finding the \emph{shortest}
one is the hard part.

Third, ``God's number'', the worst-case number of moves you would ever need, is only known exactly for the normal 3 by 3
cube: 20 moves.

Learned, AI-based solvers so far stop at 5 by 5. The best solutions for big cubes came from a Kaggle competition called
Santa 2023, which we use as our benchmark.

\slide{What we ask}{40 s}
Our talk answers four questions. \act{read the four, briefly}
One: what does it take to handle every size? Two: which method is shortest at each size? Three: can anyone beat the competition
winners? Four: how far is everyone from the true best answer?

The important word is \emph{fair}: every method gets the same puzzles, and every answer is checked by the same code.

\slide{Our contributions}{50 s}
We bring five things.
\textbf{SymPort}, our main method: it looks at one scramble from 48 different angles and keeps the shortest answer.
\textbf{NetGain}: a smarter way to fix many parts of a big cube at once. It is the best single solver on large random cubes.
\textbf{OneNet}: one neural network that works for every size, even sizes it never saw in training.
A \textbf{counting bound}: a way to say how far from perfect everyone is.
And the \textbf{open engine} that all of this runs on.
On top of that, we rebuilt five known methods and ran three outside programs through the same system.

\slide{Background: clusters}{1 min}
Two puzzles. The cube you know: each move turns one slice. The Rubik's Square is a flat, one-layer version: each move flips a
whole row or column upside down.

The key idea behind almost everything today is the \textbf{cluster}. \act{point at ``Clusters''}
A piece can only ever move among positions of its own kind: corners stay with corners, a given ring of edge pieces stays in
that ring. So the puzzle splits into many small groups that never mix. Solving the puzzle simply means solving every group.

Two ways to count moves: HTM, where any turn counts as one; and QTM, where a half turn counts as two. QTM is exactly the
competition's score.

%% ====================================================================
\section{Part 2: How the puzzles are solved, and our engine}

\slide{The two puzzles (picture)}{50 s}
\act{left picture} The Square is one cubie thick. A move flips a whole row or column over, so the top and bottom colours swap.
Doing the same flip twice undoes it.
\act{right picture} The cube: any of its slices can turn a quarter, half or three-quarter turn.
The point of this slide: they look different, but inside they have the \emph{same} cluster structure. That is why one engine
and the same ideas work for both, and why we use the small Square as a test lab where we can compute perfect answers.

\slide{Solving the Square}{1 min}
The main tool is a short move pattern: flip a column, a row, the same column, the same row. We call it ``c r c r''.
\act{point at the panel showing the net effect}
Almost everything goes back to where it was. Only three cells change places: the crossing point and its two mirror cells.
So with four moves we can make a tiny, controlled change. Tools like this are called \textbf{commutators}.

We solve from the outside in: first the outer frame, which at any size behaves like a small 3 by 3 Square and is solved
perfectly by lookup; then the edges; then each inner group with these four-move tools.

\slide{Solving the cube}{1 min}
The cube works the same way, one level up.
\act{left} These are the cluster types on a 5 by 5 cube: corners, edges, centers.
\act{middle} A \textbf{3-cycle} swaps three pieces of one cluster and touches nothing else. Our engine generates all 4,048 of them
for each cluster, automatically.
\act{right} The standard recipe, called reduction: first fix a parity problem; then solve the centers; then the cube behaves like
a normal 3 by 3, which we solve with Kociemba's well-known algorithm; finally fix the edge pieces.

\slide{Related work (optional)}{30 s}
We build on four lines of work: perfect solvers for the 3 by 3 (Korf, Kociemba, Rokicki); methods for big cubes (reduction, and
Demaine's theory method); learned solvers (DeepCubeA and CayleyPy, which stop at 5 by 5); and the Santa 2023 winners, team
KMCoders. The gap: learned solvers need one network per size, Demaine's method is theory, and the competition answers are the
best of many runs of a heavily tuned program.

\slide{The engine}{1 min}
This slide is about trust. \act{point across the six boxes}
The same code builds both puzzles at any size: a move is just a rotation of one slab of stickers.
It finds the clusters automatically, in an order that does not depend on size; that is what later lets one network work at
every size.
It generates the 3-cycles, simplifies move lists, and checks that a scramble is valid.

Most important: the \textbf{harness}. No solver is trusted. Every answer, including answers from other people's programs,
is replayed by our engine, and the engine decides whether the puzzle is really solved. A crash or a timeout counts as a failure.

\slide{Eight solvers and three baselines}{40 s}
Here is the whole cast. Five are our own re-implementations of published methods. Three are new: OneNet, NetGain and SymPort,
highlighted. Three are baselines run unchanged: the competition winners' published answers, one run of their program, and an
open-source solver called dwalton76. All of them go through the same harness.

%% ====================================================================
\section{Part 3: The methods}

\slide{IDA*-PDB}{1 min}
The first method finds the \emph{perfect} answer, but only for small puzzles.
It searches move sequences step by step, and uses a ``guess of how many moves are left'' to cut off hopeless paths early.
The guess comes from lookup tables, called pattern databases: for example, the exact number of moves needed to fix just the 8
corners, for every possible corner arrangement, which is 88 million entries. The guess never overestimates, so the first
answer found is guaranteed to be the shortest.

Results: perfect answers on every random Square up to 5 by 5, in about 13 seconds each. By searching every state, it also
gave us the exact God's numbers for Squares: 6, 10 and 12 moves for sizes 3, 4 and 5. Beyond that, the search time explodes.

\slide{Reduction}{40 s}
This is how people solve big cubes by hand. Fix parity, solve the centers, solve the cube as a 3 by 3, then fix the edges.
Each step uses 3-cycles that do not disturb what is already done.
The weakness: each 3-cycle fixes only about one piece, so the number of moves grows with the square of the size. It is
the longest of our cluster methods: 1,885 moves on a 10 by 10 cube.

\slide{Demaine}{1 min}
Demaine and colleagues had a clever idea: instead of fixing one cluster at a time, use \emph{one} move sequence that fixes the
same thing in \emph{many} clusters at once. Turn several columns together and several rows together, and one commutator acts on
a whole block of clusters. In theory this is the best possible growth rate.

On real cubes we had to adapt it: we batch identical 3-cycles across a block of clusters. But there is a rule: the block is used
only if \emph{every} cluster in it needs that fix. As solving goes on, clusters become different from each other, and the blocks
shrink to single clusters. On the 20 by 20 cube, 244 of 292 rounds fixed just one cluster. That problem is what NetGain solves.

\slide{CayleyPy}{1 min}
Now the learned methods. CayleyPy trains a neural network to guess how far a state is from solved.
Training data is free: start from the solved puzzle, make random moves, and label each state with the number of moves made.
To solve, it uses \textbf{beam search}: keep the best few thousand states according to the network, try every move from each,
keep the best again, and repeat until one is solved.

On the Square it is excellent: perfect answers at sizes 3 and 4, and wider beams help, just as the original paper reports.
The catch: every size needs its own network, and cubes need billions of training examples. So we use it on Squares only.

\slide{OneNet (new)}{1 min}
OneNet fixes that catch. Instead of feeding the network every sticker, we give it one \textbf{token per cluster}.
Because our engine lists each cluster's stickers in the same order at every size, a token means the same thing on a 5 by 5 as
on a 14 by 14. A bigger puzzle is just more tokens.

The network is a transformer without position order, because clusters are a set, not a sequence. Each token predicts its own
cost, and we add them up.

We trained it only on sizes 4 to 8, and it solves sizes it never saw, up to 14 by 14. Its distance predictions still correlate
well at unseen sizes: \textbf{0.66, 0.65 and 0.53} at sizes 10, 14 and 20. \act{if the slide still says 0.70, say these numbers}

\slide{NetGain (new)}{1 min}
Back to Demaine's problem: blocks needed \emph{every} cluster to agree. NetGain changes one rule.
We ask: does the block \emph{as a whole} pay? Add up the pieces it fixes across all clusters, and compare with the best you could
do with a single 3-cycle for the same number of moves. If the block wins, use it, even if one or two clusters gain nothing or
even lose a little.

Everything else is exactly Demaine's code, so any difference in results comes from this one rule. On the same 20 by 20 cube,
229 of 333 rounds are now bulk moves, instead of nearly all being single fixes.

\slide{NetGain in one picture}{50 s}
\act{point at the grid} Each cell is one cluster. Green cells (+) gain from this move, grey cells have no effect, and red ($-$) lose a little.
Demaine's rule only allows blocks where every cell is green; the best one here fixes 1.3 pieces per move.
NetGain allows this bigger block, with one neutral and one slightly bad cell, because overall it fixes 1.6 pieces per move.
More pieces per move means a shorter solution.

\slide{KMC-Repro}{1 min}
This is the competition winners' method, rebuilt only from their written description, before we looked at their code.
Their key trick: if you store the effect of every ending part of the move list, you can test inserting a new move sequence at
\emph{any} point with a single lookup, without replaying everything. That makes millions of edits cheap to test.
The pipeline: solve the special pieces, then use simulated annealing to improve the move list, then insert short commutators,
then finish with 3-cycles.

Our rebuild is 1.6 to 1.7 times longer than their published answers, while their real program is only 5 to 18 percent longer. The
gap shows what the write-up leaves out.

\slide{SymPort (main method)}{1.5 min}
Now our main method. \act{follow the numbers 1 to 4}
A cube has 48 symmetries: 24 ways to rotate it, times a mirror image. Take the scramble and look at it from each of the 48
views. Each view is a different-looking puzzle, but it is really the same problem.
Run the same strong solver, the winners' own program, on every view. Each gives a slightly different answer, because solvers are
sensitive to small differences in their input.
Translate each answer back to the original cube. Check every one, and keep the shortest.

Why are the lengths exactly the same after translating? \act{point at the formula}
A rotation or mirror always turns a slice move into another slice move of the same size. So a 60-move answer for a rotated
copy becomes a 60-move answer for the original. Nothing is lost.

On the cube, the inner solver is KMCoders' program with three settings: their beam version, a shorter annealing run, and a
perfect solver for the central 3 by 3 when the size is odd.

\slide{SymPort in one picture}{50 s}
A real example: Santa 4 by 4 puzzle 151. \act{point at the dots}
Each dot is one of the 48 views, solved and translated back. The lengths range from 60 to 88 moves. One normal run of their
program gave 90. Their published best, from many runs, was 68. SymPort simply keeps the 60.

\slide{Counting lower bound}{1 min}
Question four: how far from perfect is everyone?
The idea is counting. There are only so many different move sequences of length 10, or 20, or 100. Every state must be reached
by at least one of them. So if the number of short sequences is smaller than the number of states, some states must need more
moves. That gives a minimum, a lower bound, for any size.

We checked it on Squares where we know the truth. \act{point at the table}
The true worst case is 1.5 to 1.7 times this bound. And the exact values for the 4 by 4 and 5 by 5 Squares are new; we have not
seen them published before.

%% ====================================================================
\section{Part 4: Results}

\slide{Results: Rubik's Square}{1 min}
On Squares, SymPort is the best at all 18 sizes, from 3 to 20.
\act{point at the top rows}
At small sizes, up to 5, it reaches the perfect answer, as CayleyPy's networks do.
\act{point at the bottom rows}
At 9 and above no per-size network exists, and SymPort beats NetGain, the best single solver, by up to about 7 percent.
NetGain in turn beats Demaine and Reduction.

\slide{Reproductions vs originals (optional)}{40 s}
Before comparing, we checked that our rebuilt methods behave like the originals.
CayleyPy: wider beams help, as in its paper. IDA*: our table sizes match Korf's exactly, and the 3 by 3 Square's worst case
matches the published value. KMCoders: one run of their program is 5 to 18 percent longer than their published answers. Our
rebuild from the write-up is 1.6 to 1.7 times longer.

\slide{Results: Santa 2023 cubes}{1.5 min}
This is the main result. Lengths are in quarter turns, the competition's score; lower is better.
\act{point at the SymPort and KMC published columns}
For sizes 4, 5, 6 and 7, SymPort is shorter on average than the published winners' answers. For example, at 4 by 4: 64.1 against
66.7.
At 8 and 9 the published answers win, and SymPort is second, still ahead of one run of their program.
After two losses in a row, we stopped at 9. Every other method is far behind: our cluster methods are about three times longer
on these puzzles.

\slide{SymPort puzzle by puzzle}{50 s}
Over all 78 puzzles from 4 to 9, SymPort uses 8,302 moves in total against 8,452 published, 1.8 percent fewer: 43 wins, 14 ties,
21 losses.
For sizes 4 to 7 alone, 72 puzzles, it saves 2.8 percent.
At 8 and 9 it is 1 to 2 percent longer, winning one puzzle of three at each size.

\slide{Why the portfolio works}{50 s}
\act{left chart} Each column is one 4 by 4 puzzle; each dot is one of the 48 views. The answers spread over 18 to 28 moves. The bar
is the published best: it sits low, but the best dot often goes below it.
\act{right chart} If we take the best of more and more views, the average keeps falling, and it has not flattened even at 48. So
the variety between views is what wins.

\slide{Where SymPort stops winning}{1.5 min}
Why does it stop at size 8? Think of two numbers.
The \textbf{gap}: how far a typical view is behind the published answer.
The \textbf{gain}: how much picking the best view saves over a typical view.
SymPort wins only when the gain is bigger than the gap.

Up to 7 by 7, it is. At 8 and 9, the gain is 3.3 percent, but the gap is about 5 percent. Three reasons:
First, the views do not get more different as the cube grows. They differ by about 7 moves at every size, while the solutions grow
much longer, so in percent the variety shrinks from 6 percent to 2 percent.
Second, each view gets expensive: 31 to 47 minutes of computing at sizes 8 and 9, so we could afford only 12 views.
Third, the published answer is itself the best of many runs.
Even with all 48 views, our model predicts a tiny win at 8 and still a loss at 9. So symmetry alone runs out; it needs more
varied inputs or a stronger base solver.

\slide{NetGain on random cubes}{50 s}
Now random cubes. This chart shows how many fewer moves NetGain needs compared with Demaine. We tested 20 cubes per size and
measured the difference on the same cubes, with 95 percent confidence intervals.
From size 5 on, NetGain is clearly shorter, and the lead grows with size: about 17 percent at size 10. On the 20 by 20 cube: 4,925
moves against 6,612.

\slide{Scaling with N}{50 s}
Here every method's length is divided by the theoretical best growth rate. A flat line would mean a method grows at the ideal
rate.
Every line still rises, so no method reaches the ideal in practice. On random cubes NetGain bends over the most, so it comes
closest. On the Santa puzzles, SymPort and the published answers are the lowest.

\slide{How far from optimal}{50 s}
Compared with our lower bound, the published answers are 1.81 times the bound at size 4, 1.99 at 7, and 2.22 at 10.
At size 4 that is close to optimal; remember, on Squares even the perfect answer is 1.5 to 1.7 times the bound.
But the gap grows with size. The reason: the last stage of the competition method fixes only a few pieces per move. Fixing many
clusters at once, as NetGain does, is how to close it.

%% ====================================================================
\section{Part 5: Lessons, limits, next steps}

\slide{Discussion}{1 min}
Four lessons.
One: write-ups leave things out. Our rebuild from their description was far behind their real code.
Two: the published answers are already locally tight. Polishing them removed no moves; running the whole solver on different
inputs is what finds shorter ones.
Three: on big random cubes, structure wins: fixing many clusters with one move sequence.
Four: learning helps at small sizes. Per-size networks are perfect up to 5; OneNet generalises, but above 7 the hand-designed
methods are shorter.

\slide{Limitations}{50 s}
We want to be upfront.
SymPort wins only up to 7 by 7, and we stopped at 9. From size 6 on we had only 3 or 4 puzzles per size, too few to call the
wins significant.
We have not yet tested whether 48 random seeds of the same program would give the same variety as 48 symmetries. That is the most
important open test.
SymPort's core solver is KMCoders' program; we release only our wrapper.
Our bound is only a lower bound: it cannot prove any single answer is near-perfect.
And our learned solvers are for Squares only.

\slide{Future work}{40 s}
What we would try next, following the analysis:
more varied copies, and first a fair test of symmetries against random seeds;
combining the best \emph{stage} from each copy instead of only picking one whole answer;
spending the computing budget adaptively, finishing only the most promising copies;
a stronger base solver, for example NetGain's bulk moves in the final stage;
then the remaining Santa sizes and learned solvers for cubes.

\slide{Conclusion}{40 s}
One engine, eight solvers, three answers.
On random cubes, the winner at large sizes is structural: batch moves when the block pays, not only when everyone agrees.
On the competition puzzles, the cube's 48 symmetries turn the winners' own solver into a portfolio that beats their published
best up to 7 by 7.
And on optimality: little room left at size 4, but a gap that grows with size.
All our code and results are open, at the link on the slide. Thank you; we are happy to take questions.

%% ====================================================================
\newpage
\section{Quick answers to the most likely questions}
\textbf{``Isn't SymPort just running someone else's solver many times?''}
Yes, the inner solver is theirs, and we say so clearly. What is ours: the idea that symmetries are a free, seed-free source of
variety; the proof that answers translate back with exactly the same length; the configuration; and the measurement of where and
why it stops working.

\textbf{``Why the Rubik's Square at all?''}
It has the same structure as the cube but is small enough to solve perfectly up to 5 by 5. So it is where we can check against
true optimal answers and train networks cheaply.

\textbf{``How do you know the answers are correct?''}
Our engine replays every answer, including other people's, and decides itself whether the puzzle is solved.

\textbf{``Why did you stop at size 9?''}
Two losses in a row with a growing gap, and our model says even 48 copies would not win at 9. Each extra puzzle costs many
CPU-hours.

\textbf{``You reran some copies: isn't that cherry-picking?''}
Only copies that the computer killed for running out of memory were rerun. Every copy that finished was kept, good or bad. Size 8
did not improve at all.

\textbf{``Is SymPort optimal?''}
No, and we never say that. It is the best method we know of up to 7 by 7. Every method is still about twice the lower bound.

\textbf{``What is NP-complete / God's number?''}
God's number is the worst case: the fewest moves that solve every possible state; 20 for the 3 by 3. NP-complete means checking a
proposed answer is fast, but no fast method is known for finding the shortest one; finding \emph{some} answer is easy.

\end{document}
